\textbf{Scale and seeds.} Four seeds in the main grid and at $N=100,2000$, three in the sweeps and at $N=250,1000$, against five or more requested; with these $n$ the Welch tests have little power, per-seed variance under flip noise is large (std up to \rhval{analysis/analysis/flip_p0.3/ce_miou_sd_flip_p03_n1000/mean:2} mIoU), and a non-significant result is not evidence of equality. No multiplicity correction was applied.

\textbf{One architecture, one domain, one budget.} A \rhval{const/param_count}-parameter U-Net on $64\times64$ synthetic shapes with a fixed \rhval{const/train_steps}-step budget for every $N$. We did not tune the optimiser per system or per noise level, so baselines such as SCE (and the large-$q$ GCE) may be handicapped by the short budget; the optimiser budget was selected from pilot runs including two single-seed clean-label runs scored on the test split, a small leak that applies equally to all systems. A full study would use a held-out clean validation set for model selection, a longer schedule per $N$, real datasets with real annotation noise, and several architectures.

\textbf{Noise models and levels.} Boundary noise was run at one level and only at $N=100,500,2000$; the two noise models are not matched in the number of corrupted pixels, so the flip-vs-boundary comparison (H2) compares nominal levels, not equal corruption. Flips are class-symmetric and object-wide, a simple model of real annotation errors. We reimplemented only GCE, SCE and our band-ignore loss; no sample-selection, co-teaching or early-learning method and no Dice-type loss was evaluated, so the negative result for robust losses covers only the pixel-wise losses tested.

\textbf{Confounds and untested mechanisms.} Fixed steps means small $N$ is trained for more epochs, confounding size with epochs; the explanation for band-ignore CE's failure and the claim that flip damage is a foreground-class confusion beyond the per-class IoU are not isolated by any run.
